\documentclass[11pt]{article}

% --- Packages ---
\usepackage[margin=1in]{geometry}
\usepackage{amsmath, amssymb}
\usepackage{graphicx}
\usepackage{booktabs}
\usepackage{hyperref}
\usepackage{authblk}

% --- Title ---
\title{\textbf{Residue Manifold Learning: \\ 
Constraint Sampling, Structural Recovery, and Phase-Lock Geometry}}

\author{Dan Hawkley}
\affil{Independent Researcher \\ Boulder, CO}

\date{\today}

\begin{document}

\maketitle

% --- Abstract ---
\begin{abstract}
We study structure, sampling, and representation in modular arithmetic through the mod30 residue manifold. Prime-support residues form an eight-lane discrete manifold embedded in $\mathbb{Z}/30\mathbb{Z}$. We show that constraint-aligned sampling improves signal access, that nonnegative matrix factorization (NMF) recovers the manifold compactly, and that sparse autoencoder (SAE) configurations can fragment or dilute structure despite comparable reconstruction performance. We introduce the Constraint-Guided Coherence Score (CGCS), a bounded metric for structural fidelity, and demonstrate that representation regimes separate into phase-locked, partial, and diluted regions. Finally, we give a geometric interpretation of CGCS as a cosine alignment constraint with a phase-lock boundary near $45^\circ$. This connects sampling, representation, and geometry in a unified framework.
\end{abstract}

% --- Introduction ---
\section{Introduction}

Understanding how structure, sampling, and representation interact is central to modern computational learning. Recent work has shown that apparent computational advantages can arise from access to structured data representations rather than fundamentally different computational power. 

We investigate this principle in a controlled modular setting. Prime residues modulo 30 form a discrete, structured manifold. This provides a minimal environment to study:

\begin{itemize}
\item Structure: existence of constrained residue lanes
\item Sampling: access aligned vs unaligned with structure
\item Representation: recovery vs fragmentation of structure
\end{itemize}

We show that these elements interact to determine structural fidelity.

% --- Residue Manifold ---
\section{Residue Manifold Structure}

Let $n \in \mathbb{Z}$ and consider residues modulo 30:

\[
r = n \bmod 30.
\]

Prime numbers lie only in the residue set:

\[
\{1,7,11,13,17,19,23,29\}.
\]

This defines an eight-lane discrete manifold embedded in $\mathbb{Z}/30\mathbb{Z}$.

\begin{figure}[h]
\centering
\includegraphics[width=0.7\textwidth]{figures/residue_circle_mod30.svg}
\caption{Residue manifold embedded on a circle. Valid prime-support residues form eight discrete angular lanes.}
\end{figure}

% --- Sampling ---
\section{Constraint Sampling}

We compare two sampling regimes:

\begin{itemize}
\item Uniform sampling over all residues
\item Constrained sampling restricted to valid lanes
\end{itemize}

Uniform sampling wastes probability mass outside the manifold:

\[
\text{signal rate} \approx \frac{8}{30}.
\]

Constrained sampling achieves:

\[
\text{signal rate} = 1.
\]

\begin{figure}[h]
\centering
\includegraphics[width=0.75\textwidth]{figures/sampling_signal_rate.svg}
\caption{Signal rate under uniform vs constrained sampling.}
\end{figure}

% --- Representation ---
\section{Representation and Recovery}

\subsection{NMF Recovery}

Nonnegative matrix factorization (NMF) recovers the residue manifold as a compact basis:

\[
X \approx WH,
\]

where $H$ aligns with residue lanes.

\begin{figure}[h]
\centering
\includegraphics[width=0.75\textwidth]{figures/nmf_components.svg}
\caption{NMF components align with residue lanes.}
\end{figure}

\subsection{SAE Dilution}

Sparse autoencoders (SAE) may:

\begin{itemize}
\item duplicate features
\item leave features inactive
\item spread features across residues
\end{itemize}

\begin{figure}[h]
\centering
\includegraphics[width=0.75\textwidth]{figures/sae_dictionary_features.svg}
\caption{SAE dictionary features showing redundancy and diffusion.}
\end{figure}

% --- Phase Diagram ---
\section{Phase Diagram of Structural Regimes}

We combine results into a phase diagram over model capacity:

\begin{itemize}
\item Coverage of valid lanes
\item Alignment to residue structure
\item Redundancy and inactive features
\end{itemize}

\begin{figure}[h]
\centering
\includegraphics[width=0.75\textwidth]{figures/coverage_phase_diagram.svg}
\caption{Coverage vs capacity for NMF and SAE.}
\end{figure}

We observe distinct regimes:
\begin{itemize}
\item Recovered (compact, aligned)
\item Partial
\item Diluted
\item Fragmented
\end{itemize}

% --- CGCS ---
\section{Constraint-Guided Coherence Score (CGCS)}

We define:

\[
\text{CGCS} = 
(\text{coverage})
\times (\text{alignment})
\times (\text{redundancy penalty})
\times (\text{dead feature penalty})
\times (\text{reconstruction penalty})
\]

CGCS is bounded in $[0,1]$.

We define a phase-lock threshold:

\[
\text{CGCS} \ge \frac{24}{25}.
\]

\begin{figure}[h]
\centering
\includegraphics[width=0.75\textwidth]{figures/cgcs_quality_gate.svg}
\caption{CGCS separates phase-locked and degraded regimes.}
\end{figure}

% --- Geometry ---
\section{Phase-Lock Geometry}

We interpret CGCS geometrically using cosine similarity:

\[
\cos \theta = \frac{A \cdot B}{\|A\|\|B\|}.
\]

Structure quality corresponds to alignment angle $\theta$.

\begin{itemize}
\item $\theta \approx 0^\circ$ — fully aligned
\item $\theta \le 45^\circ$ — phase-locked
\item $\theta > 45^\circ$ — degraded
\end{itemize}

\begin{figure}[h]
\centering
\includegraphics[width=0.75\textwidth]{figures/phase_lock_region.svg}
\caption{Phase-lock region defined by angular threshold.}
\end{figure}

% --- Discussion ---
\section{Discussion}

We show that:

\begin{itemize}
\item Structure exists independently of representation
\item Sampling determines access to structure
\item Representation determines whether structure is preserved
\end{itemize}

Compact factorizations preserve structure, while overparameterized sparse representations may dilute it.

% --- Conclusion ---
\section{Conclusion}

Residue-manifold learning provides a minimal setting to study structural fidelity. CGCS unifies coverage, alignment, and representation quality into a single metric. Its geometric interpretation reveals a phase-lock boundary governing structure preservation.

% --- References ---
\begin{thebibliography}{9}

\bibitem{tang}
E. Tang,
\textit{A quantum-inspired classical algorithm for recommendation systems},
arXiv:1807.04271 (2018).

\end{thebibliography}

\end{document}
